\chapter*{General Conclusion}
\addcontentsline{toc}{chapter}{General Conclusion}

This report presented CANvisionNative, a hybrid-architecture intrusion detection system for vehicle \gls{can} bus networks, from its motivation through its design, implementation, and validation.

\section*{Achievements}
\addcontentsline{toc}{section}{Achievements}

The project delivered a working desktop client and backend pipeline supporting three input modes (replay, live simulation, and live hardware capture through a shared code path), a four-layer detection pipeline (timing, payload, machine learning, temporal persistence) fused into a single score, a vehicle-specific model store trained on more than seven million \gls{can} frames across eighteen profiles, and a natural-language explanation layer combining rule-based reasoning with a Large Language Model provider chain that degrades gracefully when no external model is reachable. Beyond the software itself, the project delivered a rigorous, evidence-based validation of that software against a real one-hour vehicle recording, which is documented in full in Chapter~\ref{chap:validation}.

\section*{Objectives reached}
\addcontentsline{toc}{section}{Objectives reached}

Every objective listed in Chapter~\ref{chap:host} was reached: the four-layer fusion pipeline was built and validated; recorded logs, live simulation, and live hardware capture are all supported through the same detection engine; vehicle-specific models are trained, selected at run time, and the selection was proven correct with direct evidence after a defect was found and fixed; the desktop client supports import, replay, review, and export; the explanation layer produces a plain-language description and recommendation for every alert; and the complete pipeline was validated end to end against real data, with every defect found fixed and re-verified rather than worked around.

\section*{Scientific contributions}
\addcontentsline{toc}{section}{Scientific contributions}

Chapter~\ref{chap:sota} positioned this project against the reviewed literature on three points where most anomaly-based \gls{can} \gls{ids} designs stop short: combining several independent detection signals into one fusion score instead of relying on a single signal, selecting a genuinely per-vehicle trained model at run time instead of one generic model, and producing a human-readable explanation of a detected anomaly instead of a bare numeric score. The root-cause methodology used in Chapter~\ref{chap:validation} — explaining 100\% of the alerts produced on a real recording by their dominant contributing layer before deciding whether any correction was needed — is itself a contribution that could be reused to audit other anomaly-detection systems built on the same fusion-score design.

\section*{Engineering contributions}
\addcontentsline{toc}{section}{Engineering contributions}

Six confirmed software defects were found through evidence-based testing and corrected: incorrect vehicle-model selection, driving-state misclassification, mislabeled root-cause attribution (found duplicated across four files), a replay-parser inconsistency across log formats, a permanently-stopped alert-polling timer, and a stale-data leak in the \gls{ai} explanation cache. Each was diagnosed to its exact root cause before being changed, and each fix was verified with direct, reproducible evidence — measured on the same real recording before and after — rather than by tuning detection thresholds. This produced a measured 25\% reduction in alert volume and eliminated machine-learning score saturation entirely, without artificially suppressing any alert.

\section*{Limitations}
\addcontentsline{toc}{section}{Limitations}

As detailed in Section~\ref{chap:validation}, the live hardware capture path remains unvalidated on a physical vehicle, pending the arrival of a compatible adapter, and the root-cause proportions found on the Kia EV6 recording are not guaranteed to generalize to every vehicle without further testing.

\section*{Future work}
\addcontentsline{toc}{section}{Future work}

The immediate next step is the hardware validation described in Chapter~\ref{chap:requirements} as milestone M2, followed by extending the evidence-based root-cause methodology of Chapter~\ref{chap:validation} to additional vehicle profiles, and refining the scope of the historical-case reference feature discussed in Section~\ref{sec:validation-ai}.

\bigskip
Overall, this project shows that a desktop-class, multi-layer, explainable intrusion detection system for vehicle \gls{can} bus networks can be built with commodity hardware and open datasets, and — just as importantly — that it can be held to the same standard of evidence as any other safety-relevant piece of software before it is trusted.

\cleardoublepage
